\begin{prob}
    Write each of the following functions in $\Theta$-notation in simplest terms.
    You do not need to show your work or provide justification.

    Hint: your answer should have the form $\Theta(n^\alpha)$, $\Theta((\log_\mu(n))^\beta)$,
    $\Theta(\gamma^{\delta n})$, or $\Theta(1)$, where $\alpha, \beta, \gamma, \delta, \mu$ are
    constants that depend on the problem. While you \emph{can} find the answer
    using more formal means (such as limits), you should be able to find the
    answers rather quickly using informal reasoning. You can always use the
    formal approach to check your answer, if you wish.

    \begin{subprobset}

        \begin{subprob}
            $f(n) = \sqrt{n^3 + 2n^2 + 500}$

            \begin{soln}
                $f(n) = \Theta(n^{3/2})$
            \end{soln}
        \end{subprob}

        \begin{subprob}
            $f(n) = 20 \times 2^{\log_2 (n^3 + n) }$

            \begin{soln}
                $f(n) = \Theta(n^3)$
            \end{soln}
        \end{subprob}

        \begin{subprob}
            $f(n) = 2^{10} + 10^2 + \log_{10} 1000$

            \begin{soln}
                $f(n) = \Theta(1)$.
            \end{soln}
        \end{subprob}

        \begin{subprob}
            $f(n) = \log_2(3^{n+2} + 5n^3 + 1)$

            \begin{soln}
                $f(n) = \Theta(n)$.

                $3^{n+2}$ grows a lot faster than the other terms inside of the logarithm,
                so we can ignore them. Therefore we ask: what is $\log_2(3^{n+2})$ simplified
                in $\Theta$-notation?

                If you look up a table of logarithm properties, you'll find
                the identity $\log_a b^n = n \log_a b$.  Therefore $\log_2(3^{n+2}) = (n+2) \log_2 3 = \Theta(n)$.

                In general, $\log_a b^{g(n)} = \Theta(g(n))$, as long as $a$ and $b$ are constants greater than one
                independent of $n$.
            \end{soln}
        \end{subprob}

        \begin{subprob}
            $f(n) = 2^{10n} + 10^{2n}$

            \begin{soln}
                $f(n) = \Theta(2^{10n})$

                Informally,
                $2^{10n}$ grows a lot faster than $10^{2n}$, so this is dominated by $2^{10n}$.
                Note that $\Theta(2^{10n})$ is \emph{not} the same as $\Theta(2^{n})$; the former
                grows a lot faster. This is a case where constants \emph{do} matter.

                If you want to see this more formally, the easiest way might be to use
                limits. We'll show that
                \[
                    \lim_{n \to \infty} \frac{2^{10n} + 10^{2n}}{2^{10n}} = 1
                \]
                which proves that $f(n) = \theta(2^{10n})$.

                \begin{align*}
                    \lim_{n \to \infty} \frac{2^{10n} + 10^{2n}}{2^{10n}}
                    &=
                    \lim_{n \to \infty} \left[\frac{2^{10n}}{2^{10n}} + \frac{10^{2n}}{2^{10n}}\right]
                    \\
                    &=
                    \lim_{n \to \infty} \frac{2^{10n}}{2^{10n}} + \lim_{n \to \infty} \frac{10^{2n}}{2^{10n}}
                    \\
                    &=
                    1 + \lim_{n \to \infty} \frac{10^{2n}}{2^{10n}}
                    \\
                \intertext{This limit is easier to find if we work with the same base for each exponential.
                We can write $10 = 2^{\log_2 10} \approx 2^{3.3}$.}
                    &=
                    1 + \lim_{n \to \infty} \frac{(2^{\log_2 10})^n}{2^{10n}}
                    \\
                    &=
                    1 + \lim_{n \to \infty} \frac{2^{(\log_2 10)n}}{2^{10n}}
                    \\
                    \intertext{Recall that $2^a/2^b = 2^{a - b }$, so:}
                    &=
                    1 + \lim_{n \to \infty} 2^{n(\log_2 10 - 10)}
                    \\
                    \intertext{$\log_2 10 < 10$, so the exponent is negative and the limit goes to zero:}
                    &=
                    1 + 0
                    \\
                \end{align*}
            \end{soln}
        \end{subprob}


        \begin{subprob}
            $f(n) = \log_2(n^2 + 3n + 5) \times \log_2(3n^3)$

            \begin{soln}
                $\Theta((\log_2 n)^2)$. The base doesn't actually matter, though, so
                we typically write $\Theta((\log n)^2)$.
            \end{soln}
        \end{subprob}

    \end{subprobset}
\end{prob}
